\section*{Capítulo \ref{ch:b2:reproductive-hormones} --- Control hormonal de la reproducción}

\begin{solution}{exo:b2:reproductive-hormones:1}
Hipotálamo (GnRH, pulsos) $\to$ adenohipófisis (LH, FSH) $\to$
\omterm{def:b2:mammal-reproduction:testis}{testículo}: la LH sobre las \omterm{def:b2:mammal-reproduction:testis}{células de Leydig} (testosterona), la FSH sobre
las \omterm{def:b2:mammal-reproduction:testis}{células de Sertoli} (\omterm{def:b2:mammal-reproduction:testis}{espermatogénesis}, inhibina). Retroalimentación: la
testosterona inhibe la GnRH y la LH; la inhibina inhibe la FSH.
\end{solution}

\begin{solution}{exo:b2:reproductive-hormones:2}
Ovario: \omterm{def:b2:reproductive-hormones:cycle}{fase folicular}, días 1--14, estradiol del \omterm{def:b2:mammal-reproduction:ovary}{folículo} en
crecimiento; \omterm{def:b2:mammal-reproduction:ovary}{ovulación}, día 14, \omterm{def:b2:reproductive-hormones:cycle}{pico de LH}; \omterm{def:b2:reproductive-hormones:cycle}{fase lútea}, días 15--28,
progesterona del \omterm{def:b2:mammal-reproduction:ovary}{cuerpo lúteo}. Útero: menstruación, días 1--5 (retirada
de los esteroides); fase proliferativa, días 6--14 (estradiol); fase
secretora, días 15--28 (progesterona).
\end{solution}

\begin{solution}{exo:b2:reproductive-hormones:3}
El estrógeno y el progestágeno diarios mantienen bajas la LH y la FSH por
retroalimentación negativa, de modo que no se recluta ningún \omterm{def:b2:mammal-reproduction:ovary}{folículo}, no
se produce la subida del estradiol y el pico nunca se dispara; el
progestágeno espesa además el moco. En la semana sin píldoras los
esteroides caen y el \omterm{def:b2:reproductive-hormones:cycle}{endometrio}, construido bajo su acción, se desprende
--- un sangrado por deprivación, no una menstruación tras una \omterm{def:b2:mammal-reproduction:ovary}{ovulación}.
\end{solution}

\begin{solution}{exo:b2:reproductive-hormones:4}
La LH y la FSH se multiplican por varias veces (sin retroalimentación de
la testosterona ni de la inhibina); el músculo, la libido y las glándulas
accesorias regresan. La testosterona restablece los caracteres y hace
bajar de nuevo la LH, pero no la fertilidad: los \omterm{def:b2:mammal-reproduction:testis}{testículos} ya no están,
e incluso en un hombre con \omterm{def:b2:mammal-reproduction:testis}{testículos} la testosterona exógena suprime la
LH y, con ella, la testosterona intratesticular que necesita la
\omterm{def:b2:mammal-reproduction:testis}{espermatogénesis}.
\end{solution}

\begin{solution}{exo:b2:reproductive-hormones:5}
$24/1.5 = 16$ pulsos al día; $24/4 = 6$. Cuando muere el \omterm{def:b2:mammal-reproduction:ovary}{cuerpo lúteo} los
pulsos siguen siendo lentos, lo que favorece la FSH: la FSH sube primero
--- la hormona adecuada, puesto que es la FSH la que recluta la siguiente
cohorte de \omterm{def:b2:mammal-reproduction:ovary}{folículos}.
\end{solution}

\begin{solution}{exo:b2:reproductive-hormones:6}
\qty{200}{pg/mL} $= \qty{200}{ng/L}$: $200\times 10^{-9}/272 =
\qty{7.4e-10}{mol/L} = \qty{740}{pmol/L}$. \qty{15}{ng/mL} $=
\qty{15}{\micro g/L}$: $15\times 10^{-6}/314 = \qty{48}{nmol/L}$. En
un microlitro: $7.4\times 10^{-16}$ mol $= 4.4\times 10^{8}$
moléculas.
\end{solution}

\begin{solution}{exo:b2:reproductive-hormones:7}
\omterm{def:b2:mammal-reproduction:ovary}{Ovulación} el día $35 - 14 = 21$; fértil del día 18 al día 22. Con ciclos
irregulares, la \omterm{def:b2:reproductive-hormones:cycle}{fase folicular}, y por tanto el día de la \omterm{def:b2:mammal-reproduction:ovary}{ovulación}, no
puede predecirse a partir de los ciclos anteriores, y la ventana se pierde
o se sitúa mal.
\end{solution}

\begin{solution}{exo:b2:reproductive-hormones:8}
Del día 21 al 26 hay 2,5 duplicaciones: $2^{2.5} = 5.7$ IU/L --- suficiente.
Una \omterm{def:b2:mammal-reproduction:implantation}{implantación} el día 25 da $1.4$ IU/L el día 26; el \omterm{def:b2:mammal-reproduction:ovary}{cuerpo lúteo}, que ya
está regresando, no es rescatado, la progesterona cae y el embrión se
pierde con la regla. La \omterm{def:b2:mammal-reproduction:implantation}{implantación} tardía es una causa frecuente de
pérdida precoz.
\end{solution}

\begin{solution}{exo:b2:reproductive-hormones:9}
La hipófisis responde al patrón de la GnRH, no a su cantidad: los pulsos
horarios mantienen la LH y la FSH, y la misma dosis diaria en infusión
continua las suprime. Descartó una simple relación dosis--respuesta.
Un mes de agonista de acción prolongada: tras unos días de estimulación
inicial, la LH cae casi a cero (desensibilización de los receptores), el
estradiol a niveles posmenopáusicos, el \omterm{def:b2:reproductive-hormones:cycle}{endometrio} se atrofia y no hay
ciclo --- una menopausia médica reversible.
\end{solution}

\begin{solution}{exo:b2:reproductive-hormones:10}
Sin hormona: $R^* = 100/0.5 = 200$. Continua: $100/6.5 = 15$. Un pulso de
dos minutos retira $6\times 200\times(2/60) = 40$ receptores
(\qty{20}{\%}); en la hora siguiente el déficit se reduce en un factor
$e^{-0.5}$, es decir, se recupera el \qty{39}{\%}. La reserva se
estabiliza donde la pérdida por pulso iguala la recuperación por hora:
unos 150 receptores antes de cada pulso, tres cuartas partes del máximo
--- diez veces el nivel con hormona continua.
\end{solution}

\begin{solution}{exo:b2:reproductive-hormones:11}
Menopausia: sin \omterm{def:b2:mammal-reproduction:ovary}{folículos} no hay estradiol ni inhibina, y por tanto no hay
retroalimentación: la LH y la FSH están altas. Prolactinoma: la
prolactina ralentiza el generador de pulsos de GnRH, de modo que la LH es
baja, los \omterm{def:b2:mammal-reproduction:ovary}{folículos} no son estimulados, el estradiol es bajo y no hay
ciclo. Tumor de células de la granulosa: el estradiol alto y constante
suprime la LH por retroalimentación negativa (y sin un \omterm{def:b2:mammal-reproduction:ovary}{folículo} en
maduración no hay nada que ovular), de modo que el ciclo se detiene con el
estradiol alto --- y el \omterm{def:b2:reproductive-hormones:cycle}{endometrio} prolifera en exceso.
\end{solution}

\begin{solution}{exo:b2:reproductive-hormones:12}
El estradiol bajo o breve reduce la GnRH y la LH; el estradiol alto y
sostenido (por encima de \qty{200}{pg/mL} durante más de 36 horas)
cambia la respuesta del sistema kisspeptina--GnRH y de la hipófisis, que
entretanto ha almacenado LH, de modo que la misma hormona provoca ahora
una descarga. La \omterm{def:b2:mammal-reproduction:ovary}{ovulación} debe ser de todo o nada y estar sincronizada:
el óvulo debe liberarse una sola vez, un día concreto, cuando un \omterm{def:b2:mammal-reproduction:ovary}{folículo}
está maduro. Un regulador gradual (LH proporcional al estradiol) daría una
luteinización progresiva y parcial, \omterm{def:b2:mammal-reproduction:ovary}{folículos} que ovularían a medio
madurar o varios a la vez, y ningún día fértil definido.
\end{solution}

\begin{solution}{pb:b2:reproductive-hormones:1}
\textbf{1.} Folicular, días 1--14: estradiol en aumento, progesterona
baja. Lútea, días 15--28: progesterona alta.
\textbf{2.} El pico alcanzó su máximo el día 14 (empezó el día 13); la
\omterm{def:b2:mammal-reproduction:ovary}{ovulación}, unas 36 horas después del inicio, el día 15.
\textbf{3.} Días 15 a 28: 14 días, la vida del \omterm{def:b2:mammal-reproduction:ovary}{cuerpo lúteo}.
\textbf{4.} La progesterona eleva el punto de consigna termorregulador
entre \qty{0.3}{\celsius} y \qty{0.4}{\celsius}; la subida aparece uno o
dos días después de la \omterm{def:b2:mammal-reproduction:ovary}{ovulación}, de modo que confirma que la \omterm{def:b2:mammal-reproduction:ovary}{ovulación}
se ha producido, pero no puede anunciarla.
\textbf{5.} Días 11 a 14 (210, 260, 310, 220): de tres a cuatro días, más
de las 36 horas que requiere el interruptor.
\textbf{6.} Días 12 a 16.
\textbf{7.} $310\times 10^{-9}/272 = \qty{1.14e-9}{mol/L} =
\qty{1140}{pmol/L}$; $16\times 10^{-6}/314 = \qty{51}{nmol/L}$.
\textbf{8.} $65/8 = 8$. Con una semivida de una hora, la LH se reduce a la
mitad cada hora una vez que cesa la secreción, de modo que un día después
ha vuelto casi al nivel basal.
\textbf{9.} \omterm{def:b2:mammal-reproduction:ovary}{Ovulación} hacia el día 22, ciclo de unos 36 días.
\textbf{10.} \omterm{def:b2:mammal-reproduction:implantation}{Implantación} el día 22; hacia el día 27 (cinco días, 2,5
duplicaciones), unas \qty{5.7}{IU/L}.
\textbf{11.} Una progesterona que cae a \qty{2}{ng/mL} el día 27 indica que
el \omterm{def:b2:mammal-reproduction:ovary}{cuerpo lúteo} regresa según lo previsto: nada lo ha rescatado.
\textbf{12.} Progesterona en torno a \qty{20}{ng/mL} y en aumento,
estradiol en torno a \qty{200}{pg/mL}.
\textbf{13.} Al final del ciclo el estradiol, la progesterona y la inhibina
caen a la vez y la FSH escapa a su retroalimentación; a medida que la
cohorte crece, el estradiol y la inhibina suben y la FSH baja, y solo el
\omterm{def:b2:mammal-reproduction:ovary}{folículo} con más receptores de FSH (con más células de la granulosa)
sigue creciendo con una FSH en descenso --- los demás mueren.
\textbf{14.} $R^* = 200$ sin hormona; $100/6.5 = 15$ con hormona
continua.
\textbf{15.} $6\times 200\times 2/60 = 40$ receptores, el \qty{20}{\%};
constante de tiempo $1/d = \qty{2}{h}$, de modo que en una hora se
recupera el \qty{39}{\%} del déficit; la reserva se estabiliza en torno a
150 receptores, tres cuartas partes del máximo.
\textbf{16.} Los pulsos mantienen la hipófisis en unos 150 receptores, y
esta responde a cada pulso con LH; una infusión continua lleva la reserva
a 15 y las mismas células enmudecen, sea cual sea la dosis; los pulsos
restablecen la reserva y la respuesta.
\textbf{17.} Primer día: la LH y la testosterona suben (el agonista
estimula unos receptores que todavía están ahí --- la estimulación
inicial). Al cabo de un mes: receptores desensibilizados, LH casi nula y
testosterona en niveles de castración, que es el objetivo terapéutico.
\textbf{18.} La FSH sube (sin inhibina) y la LH se mantiene normal; en cada
ciclo sobreviven más \omterm{def:b2:mammal-reproduction:ovary}{folículos} a la selección, con \omterm{def:b2:mammal-reproduction:ovary}{ovulaciones} múltiples y
gemelos dicigóticos.
\textbf{19.} En un ciclo natural la caída de la FSH mata a todos los
\omterm{def:b2:mammal-reproduction:ovary}{folículos} salvo al dominante; un aporte constante de FSH mantiene toda la
cohorte creciendo hasta la madurez.
\textbf{20.} Diez \omterm{def:b2:mammal-reproduction:ovary}{folículos} producen muy pronto mucho más de
\qty{200}{pg/mL} de estradiol, lo que desencadenaría un \omterm{def:b2:reproductive-hormones:cycle}{pico de LH}
prematuro y una \omterm{def:b2:mammal-reproduction:ovary}{ovulación} antes de la recogida; el antagonista bloquea la
GnRH y no puede producirse ningún pico.
\textbf{21.} En la tabla el pico alcanzó su máximo el día 14 y la
\omterm{def:b2:mammal-reproduction:ovary}{ovulación} se produjo unas 36 horas después de su inicio; la hCG sustituye
al pico, y la recogida a las 35 horas capta los ovocitos maduros pero
todavía dentro del \omterm{def:b2:mammal-reproduction:ovary}{folículo}.
\textbf{22.} $10\times 0.7 = 7$ embriones, $7\times 0.4 = 2.8$
\omterm{def:b2:mammal-reproduction:implantation}{blastocistos}; $1 - 0.65^{2} = 0.58$.
\textbf{23.} El estrógeno y el progestágeno constantes mantienen bajas la
FSH y la LH: no crece ningún \omterm{def:b2:mammal-reproduction:ovary}{folículo}, el estradiol nunca sube hasta el
umbral, no hay pico, ni \omterm{def:b2:mammal-reproduction:ovary}{cuerpo lúteo}, ni subida de la progesterona.
\textbf{24.} Debe administrarse testosterona para mantener los caracteres
secundarios y la libido; pero no restablece la elevada concentración
intratesticular que aportaban las \omterm{def:b2:mammal-reproduction:testis}{células de Leydig}, de modo que la
\omterm{def:b2:mammal-reproduction:testis}{espermatogénesis} sigue suprimida --- que es precisamente lo que se busca.
\textbf{25.} \omterm{def:b2:mammal-reproduction:ovary}{Ovulación} el día 15; \omterm{def:b2:reproductive-hormones:cycle}{fase lútea} de 14 días; ventana fértil,
días 12--16; pico desencadenado por un estradiol por encima de
\qty{200}{pg/mL} durante más de 36 horas, condición cumplida los días
11--14.
\end{solution}
